% =========================================================================
% secciones_es/sec01_introduction.tex — Sección I, edición en español
% (traducción rigurosa de la edición inglesa canónica).
% =========================================================================
\section{Introducción}
\label{sec:intro}

Los sistemas interconectados insulares ocupan un punto extremo del espacio de diseño de la ingeniería eléctrica. Una isla no puede importar inercia, soporte reactivo ni transferencia de emergencia desde vecinos: cada megavatio de reserva, cada voltamperio reactivo de compensación y cada megahercio de capacidad de amortiguamiento deben producirse, despacharse y defenderse dentro de la misma huella sincrónica que los consume. La República Dominicana es un caso canónico: una única isla sincrónica de unos once millones de habitantes, sin interconexión submarina ni terrestre, con más del cuarenta por ciento de la carga nacional concentrada en la región metropolitana de Ozama, un ancla de inercia dominada por la central a carbón de Punta Catalina y un despliegue renovable cuya participación solar instantánea ha alcanzado el 36.6\% de la demanda en telemetría operativa \cite{ocseni}. Las consecuencias nacionales de perder el equilibrio generación--demanda no son hipotéticas: las interrupciones de servicio a escala nacional registradas en noviembre de 2025 y febrero de 2026, documentadas en los registros de eventos del operador \cite{ocseni} y en el archivo hemerográfico nacional \cite{listin2025}, movilizaron la respuesta de emergencia a la escala de un territorio completo y renovaron el debate de política pública sobre estabilidad de frecuencia, capacidad de arranque en negro y observabilidad del patrimonio de transmisión \cite{blackout2004}. La física de esta fragilidad es ya clásica \cite{kundur1994} --- el declive de la inercia del sistema degrada la velocidad de variación de la frecuencia y estrecha la envolvente de control factible de las redes dominadas por convertidores \cite{tielens2016}, \cite{milano2018}, \cite{dorfler2013}, \cite{hatziargyriou2021} --- pero los instrumentos empíricos disponibles para estudiarla siguen siendo sorprendentemente primitivos.

La infraestructura de benchmarking del sector es, en su mayoría, sintética. Los casos de prueba canónicos que todavía gobiernan la formación doctoral y los estudios comparativos --- la red WSCC de 9 barras y los sistemas IEEE de 14, 30, 57 y 118 barras --- descienden de abstracciones de mesa de dibujo de principios de los años sesenta: las reactancias son números académicos redondos, las cargas son triángulos fijos de potencia, la geografía está ausente, el clima está ausente, y nada en el conjunto de datos se ha disparado, mantenido, pagado o recibido un huracán. En el polo opuesto, los casos de red de calidad de servicio se conservan en formatos propietarios de plataformas comerciales de simulación, están gravados por la confidencialidad y rara vez se publican en formas que permitan la re-ejecución independiente. El movimiento de sistemas abiertos ha respondido con infraestructura valiosa --- la colección de casos de MATPOWER \cite{matpower2011} y la Power Grid Library de instancias de benchmarking de AC-OPF \cite{pglib2019}, entre otras --- pero sus instancias son, casi sin excepción, \emph{sintéticas o anonimizadas}: son redes que existen para ser difíciles, no redes que existen. Entre el juguete y el secreto casi no hay nada: un modelo abierto, completo y medido de un sistema real de transmisión nacional, con procedencia suficiente para la verificación independiente, es un bien público del que hoy carece la comunidad investigadora en el Caribe insular --- y, con pocas excepciones, en el mundo insular en general.

Este artículo cierra esa brecha con \textbf{SENI-142}, un benchmark abierto construido sobre tres compromisos. Primero, \textbf{fidelidad a una realidad medida}: el modelo de red es una reconstrucción calibrada y georreferenciada del unifilar CAD, derivada del patrimonio nacional de transporte troncal y de los levantamientos unifilares, y se valida contra exactamente 520 semanas completas consecutivas de 168 horas --- 87,360 estados horarios sobre 3,640 días de lunes a domingo, extraídos sin una sola hora interpolada del archivo operativo 2016--2025 del operador del sistema --- con despacho archivado unidad por unidad y hora a hora, de modo que cada estado simulado responde a un estado operativo archivado. Segundo, \textbf{honestidad estructural sobre la escala}: se resuelve el sistema nacional completo --- todas las barras de transporte, todas las líneas troncales, todos los transformadores --- no un distrito reducido del mismo, y el método de solución se formula de modo que su exactitud y su complejidad son teoremas y no aspiraciones: una formulación jerárquica de frontera-coordenada (HBCF) en la tradición diakóptica \cite{kron1963}, \cite{happ1971}, condensada por el complemento de Schur \cite{zhang2005} y factorizada sobre una partición que sigue las diez regiones de planificación territorial del SNIT instituidas por el estatuto territorial nacional, la Ley 368-22, en estricta compatibilidad con la Ley General de Electricidad 125-01 \cite{ley368}. Tercero, \textbf{reproducibilidad determinista}: cada artefacto pasa por un ciclo de vida de gemelo digital de cuatro fases con una cascada de compuertas \emph{fail-closed}, sellado SHA-512 desacoplado y promoción atómica, de modo que cualquier cifra reportada en este artículo puede regenerarse desde un estado de software bit a bit idéntico y criptográficamente atestiguado.

Las contribuciones son cinco, y cada una es falsable:
\begin{itemize}
\item \textbf{C1 --- Benchmark.} Un modelo abierto, georreferenciado y consistente con estándares de la red completa de transmisión de la República Dominicana: 142 barras nodales derivadas de 112 recintos físicos de subestaciones cercados --- desglosadas en 115 barras de transporte troncal (6 a 345~kV, 3 a 230~kV, 106 a 138~kV) y 27 puertos de frontera condensados desde la capa de 69~kV --- 135 líneas de transmisión, 58 transformadores de potencia, estratos de 345, 230 y 138~kV, con la capa de subtransmisión de 69~kV reducida mediante equivalentes que preservan las inyecciones; publicado con topología, parámetros, metadatos de procedencia y arnés de validación mediante IEEE DataPort y GitHub.
\item \textbf{C2 --- Formulación.} La partición diakóptica HBCF de una red insular completa en diez microrredes territoriales soberanas --- las regiones de planificación de la Ley 368-22 --- agrupadas en cuatro macrodominios eléctricos con un conjunto de frontera mínimo, y la demostración de que su condensación de Schur es exacta, con complejidad $O(\sum_k N_k^3 + N_B^3)$ frente a la $O(N^3)$ de la factorización densa monolítica.
\item \textbf{C3 --- Gemelo digital determinista.} Un ciclo de vida de cuatro fases (enclave hermético de espacios de nombres POSIX con una cascada \emph{fail-closed} de catorce compuertas; verificación estática post-evento $N-1$ contra un catálogo de 215 contingencias en banco de ensayo; Golden Master con sellado SHA-512 desacoplado; promoción atómica del núcleo vía \texttt{renameat2} con regresión nula) que convierte la telemetría del solver en evidencia con cadena de custodia.
\item \textbf{C4 --- Flexibilidad multivector.} El acoplamiento del hidrógeno power-to-gas (electrólisis, almacenamiento a 700~bar, re-inyección por celda de combustible), 764~MW de bombeo hidroeléctrico reversible y biometanación basada en ADM1 a un AC-OPF multitemporal, enunciado con su sistema completo de Karush--Kuhn--Tucker (KKT) y validado por una prueba de balance zonal de CAPEX/OPEX.
\item \textbf{C5 --- Rendimiento medido en silicio.} Resolución de cada estado horario nacional de AC en 2.65~ms con 473~MiB de memoria residente y residuo fasorial de frontera $\delta = 2.788 \times 10^{-14}$~pu sobre hardware de consumo (medición experimental de referencia del motor de investigación C++/Sparse del autor en su estación Linux de 8 hilos, conservada como línea base de comparación para solvers de la comunidad; el arnés abierto liberado reproduce determinísticamente los 41 controles matemáticos y estructurales, no esta cifra de tiempo de muro) --- dos a tres órdenes de magnitud dentro del ciclo de estimación de estado de 1--5~s de un sistema de gestión de energía (EMS), adecuado para el cribado estático de contingencias --- bajo un protocolo de validación por \emph{bootstrap} de bloques de 24 horas.
\end{itemize}

El resto del artículo se organiza como sigue. La Sección~\ref{sec:methodology} enuncia la metodología: la formulación invariante (modelo de grafo, reducción de 69~kV, partición HBCF, condensación de Schur, protocolo de datos, ciclo de vida determinista, regularización de límites reactivos, AC-OPF multitemporal) y la cláusula territorial que instancia el caso dominicano bajo la Ley 368-22. La Sección~\ref{sec:topology} documenta la topología del benchmark, desde el patrimonio físico de 112 subestaciones hasta el modelo nodal de 142 barras y las diez microrredes territoriales. La Sección~\ref{sec:multivector} desarrolla la constelación de flexibilidad multivector y su validación económica zonal. La Sección~\ref{sec:benchmarking} reporta rendimiento en silicio, exactitud y resiliencia a contingencias. La Sección~\ref{sec:discussion} discute implicaciones de política y limitaciones, y la Sección~\ref{sec:conclusion} concluye con el plan de liberación abierta y cuatro rutas de investigación.
